% Part: incompleteness
% Chapter: arithmetization-syntax
% Section: coding-terms

\documentclass[../../../include/open-logic-section]{subfiles}

\begin{document}

\olfileid{inc}{art}{trm}
\olsection{ٹرماں نوں کوڈ دینا}

\begin{explain}
ٹرم علامتاں دا اک خاص قسم دا سلسلہ اے۔ ایہ ٹرماں دی تشکیل
دے قاعدیاں مطابق ثابتہ نشاناں تے متغیراں توں استقرائی ڈھنگ
نال بندا اے۔ علامتاں دے سلسلیاں نوں عدداں دے روپ وچ کوڈ کیتا
جا سکدا اے۔ ایس لئی علامتاں نوں کوڈ دین دا طریقہ تے عدداں دے
سلسلیاں نوں کوڈ دین دا طریقہ اکٹھے ورتدے آں۔ ایس کرکے ٹرماں
نوں گوڈل نمبر دینا اوکھا نہیں۔ اصل اوکھ ایہ وکھاؤن وچ اے
کہ کسے عدد دا صحیح طرح بنے ٹرم دا گوڈل نمبر ہون والی خاصیت
قابلِ حساب اے، بلکہ بنیادی بازگشتی اے۔
\end{explain}

!!^{variable}s تے !!{constant}s سبھ توں سادے ٹرم نیں۔ ایہ پرکھنا
سوکھا اے کہ $x$ ایسے کسے ٹرم دا گوڈل نمبر اے یا نہیں۔
$\fn{Var}(x)$ اُدوں قائم اے جدوں $x$، $\Gn{\Obj v_i}$ ہووے،
کسے~$i$ لئی۔ دوجے لفظاں وچ، $x$ لمبائی~$1$ والے سلسلے دا کوڈ
اے تے اوس سلسلے دا اکّو رکن $(x)_0$ کسے !!{variable}~$\Obj v_i$
دا کوڈ اے۔ یعنی $x$، $\tuple{\tuple{1, i}}$ ہووے، کسے~$i$ لئی۔
ایویں $\fn{Const}(x)$ اُدوں قائم اے جدوں $x$، $\Gn{\Obj c_i}$
ہووے، کسے~$i$ لئی۔ ایہ دونوں تعلق بنیادی بازگشتی نیں، کیوں جے
ایسا کوئی $i$ موجود ہووے تاں اوہ لازماً $< x$ اے:
\begin{align*}
  \fn{Var}(x) & \defiff \bexists{i<x}{x = \tuple{\tuple{1, i}}}\\
  \fn{Const}(x) & \defiff \bexists{i<x}{x = \tuple{\tuple{2, i}}}
\end{align*}

\begin{prop}
\ollabel{prop:term-primrec}
تعلق $\fn{Term}(x)$ تے $\fn{ClTerm}(x)$ بنیادی بازگشتی نیں۔
ایہ بالترتیب اُدوں تے صرف اُدوں قائم ہوندے نیں جدوں $x$
کسے ٹرم یا بند ٹرم دا گوڈل نمبر ہووے۔
\end{prop}

\begin{proof}
علامتاں دا سلسلہ~$s$ اُدوں تے صرف اُدوں ٹرم اے جدوں ٹرماں دا
کوئی سلسلہ~$s_0$، \dots، $s_{k-1} = s$ موجود ہووے جیہڑا وکھاوے
کہ ٹرماں دی تشکیل دے قاعدیاں مطابق !!{constant}s تے !!{variable}s
توں ٹرم~$s$ کِویں بنیا۔ تشکیل دا کوئی دِتّا سلسلہ قاعدیاں مطابق ہون لئی ہر $i < k$ لئی
ایہناں وچوں کوئی اک شرط پوری ہونی چاہیدی اے:
\begin{enumerate}
\item $s_i$ کوئی !!a{variable}~$\Obj v_j$ ہووے، یا
\item $s_i$ کوئی !!a{constant}~$\Obj c_j$ ہووے، یا
\item $s_i$، $n$ ٹرماں $t_1$، \dots، $t_n$ توں بنیا ہووے۔
  ایہ ٹرم تھاں~$i$ توں پہلاں آئے ہون تے انہاں اُتے اک
  $n$-دلیلی !!{function}~$\Obj f^n_j$ لایا گیا ہووے۔
\end{enumerate}
ایہ وکھاؤن لئی کہ گوڈل نمبراں اُتے متعلقہ تعلق بنیادی بازگشتی
اے، سانوں ایس شرط نوں بنیادی بازگشتی فنکشناں، تعلقاں تے
محدود مسوّریت نال بیان کرنا پینا اے۔

فرض کرو $y$ اوہ عدد اے جیہڑا سلسلے $s_0$، \dots، $s_{k-1}$
دا کوڈ اے، یعنی $y = \tuple{\Gn{s_0}, \dots, \Gn{s_{k-1}}}$۔
ایہ گوڈل نمبر~$x$ والے ٹرم دی تشکیل دے سلسلے دا کوڈ اُدوں
تے صرف اُدوں اے جدوں ہر $i < k$ لئی:
\begin{enumerate}
\item $\fn{Var}((y)_i)$ قائم ہووے، یا
\item $\fn{Const}((y)_i)$ قائم ہووے، یا
\item کوئی $n$ تے کوئی عدد~$z = \tuple{z_1, \dots, z_n}$ ایسے
  ہون کہ ہر $z_l$ کسے $(y)_{i'}$ دے برابر ہووے، جتھے
  $i' < i$، تے
\[
(y)_i = \Gn{\Obj f^n_j(} \concat \fn{flatten}(z) \concat \Gn{)},
\]
\end{enumerate}
نال ایہ وی ضروری اے کہ $(y)_{k-1} = x$ ہووے۔
(فنکشن $\fn{flatten}(z)$ سلسلے $\tuple{\Gn{t_1}, \dots, \Gn{t_n}}$
نوں $\Gn{t_1, \dots, t_n}$ وچ بدلدا اے تے بنیادی بازگشتی اے۔)

شق (3) وچ اشاریے $j$، $n$، گوڈل نمبر $z_l$ جیہڑے ٹرماں $t_l$
دے نیں، تے کوڈ~$z$ جیہڑا سلسلے~$\tuple{z_1, \dots, z_n}$ دا اے،
ایہ سارے~$y$ توں چھوٹے نیں۔ اُتلے $k$ دی تھاں
$\len{y}$ رکھ سکدے آں۔ ایس لئی «$y$، گوڈل نمبر~$x$ والے ٹرم دی
تشکیل دے سلسلے دا کوڈ اے» والی شرط ایس طرح بیان کر سکدے آں کہ
تعلق دا بنیادی بازگشتی ہونا ظاہر ہو جاوے۔

ہن ایہ وکھانا رہ گیا اے کہ $y$ لئی بنیادی بازگشتی حد موجود اے۔
جے $x$ کسے ٹرم دا گوڈل نمبر اے، تاں اوہدی تشکیل دا ایسا
سلسلہ ضرور ہووے گا جیہدے وچ ودھ توں ودھ $\len{x}$ ٹرم ہون گے۔
ایس لئی کہ~$s$ دی تشکیل دے سلسلے دا ہر ٹرم~$s$ وچ کسے تھاں
توں شروع ہوندا اے تے دو ذیلی ٹرم اکّو تھاں توں شروع نہیں ہو سکدے۔
ظاہر اے کہ $s$ دے ہر ذیلی ٹرم دا گوڈل نمبر $\le x$ اے۔ ایس لئی ہمیشہ
ایسا تشکیل دا سلسلہ موجود اے جیہدا کوڈ
$\le p_{k-1}^{k(x+1)}$ ہووے، جتھے $k=\len{x}$ اے۔

$\fn{ClTerm}$ لئی !!{variable}s والی شق چھڈ دیو۔
\end{proof}

\begin{prob}
وکھاؤ کہ فنکشن $\fn{flatten}(z)$، جیہڑا سلسلے
$\tuple{\Gn{t_1}, \dots, \Gn{t_n}}$ نوں $\Gn{t_1, \dots, t_n}$
وچ بدلدا اے، بنیادی بازگشتی اے۔
\end{prob}

\begin{prop}\ollabel{prop:num-primrec}
  فنکشن $\fn{num}(n) = \Gn{\num{n}}$
  بنیادی بازگشتی اے۔
\end{prop}

\begin{proof}
  اسیں $\fn{num}(n)$ دی تعریف بنیادی بازگشت نال کردے آں:
  \begin{align*}
    \fn{num}(0) & = \Gn{\Obj 0}\\
    \fn{num}(n+1) & = \Gn{\prime(} \concat \fn{num}(n) \concat \Gn{)}.
  \end{align*}
\end{proof}

\end{document}
